\uuid{qisJ}
\titre{Matrice d'une application linéaire}
\auteur{hanani ; mzari}
\organisation{exo7}
\licence{CC BY-NC-SA 4.0}
\datecreate{2020-01-20}
\niveau{L1}
\module{Algèbre}
\chapitre{Matrice}
\sousChapitre{Matrice et application linéaire}
\difficulte{2}
\theme{Matrice d'une application linéaire}
\exo7qcmid{lille-2.matrice-application.011}

\contenu{
\begin{question}
\qclassification{108.03}
\qtags{motcle=Matrice d'une application linéaire}
On considère $\Rr^3$ muni de sa base canonique notée ${\cal {B}}$ et
$f$ l'application linéaire définie  par :
$$\begin{array}{rccc}f:&\Rr^3 &\to& \Rr^3\\
& (x,y,z)&\to &(y+z,x+z,x+y).  \end{array}$$
Soit ${\cal {B'}} = \{u_1,u_2,u_3\}$, où
$ u_1=(1,0,0), u_2=(1,1,0), u_3=(1,1,1)$. Quelles sont les assertions vraies ?
\begin{answers}
\good{${\cal {B'}}$ est une base de  $\Rr^3$}.
\bad{La matrice de $f$ dans la base $ {\cal {B}}$ est : $
\left(\begin{array}{rcc}
0&1&1\\
1&0&1\\
1&1&1\\
\end{array}\right)$}.
\good{La matrice de $f$ de la base $ {\cal {B'}}$ dans la base $ {\cal {B}}$ est : $
\left(\begin{array}{rcc}
0&1&2\\
1&1&2\\
1&2&2\\
\end{array}\right)$}.
\bad{La matrice de $f$ dans la base $ {\cal {B'}}$ est : $
\left(\begin{array}{rcc}
-1&0&0\\
0&-1&0\\
1&2&1\\
\end{array}\right)$}.
\end{answers}
\begin{explanations} La matrice de $f$  d'une base $ {\cal {B}}=(u_j)$ dans une autre base $ {\cal {B}}'=(v_i)$ est la matrice  dont la $j$ième colonne est constituée des coordonnées de  $f(u_j)$ dans la base $ {\cal {B}}'$. On déduit que :
\vskip0mm
La matrice de $f$ dans la base canonique est : $
\left(\begin{array}{rcc}
0&1&1\\
1&0&1\\
1&1&0\\
\end{array}\right)$.
\vskip0mm
La matrice de $f$ de la base $ {\cal {B'}}$ dans la base $ {\cal {B}}$ est : $
\left(\begin{array}{rcc}
0&1&2\\
1&1&2\\
1&2&2\\
\end{array}\right).$
\vskip0mm
La matrice de $f$ dans la base $ {\cal {B'}}$ est : $
\left(\begin{array}{rcc}
-1&0&0\\
0&-1&0\\
1&2&2\\
\end{array}\right)$.
\end{explanations}
\end{question}
}
